% !TeX root = thesis.tex

\chapter{Conclusion}
\label{ch:conclusion}


\section{Summary of Contributions}
\label{sec:conclusion_contributions}

This report presents an automated end-to-end framework for generating
modality-grounded question--answer (QA) annotations from multi-sensor streams
and establishing controlled Vision-Language Model (VLM) benchmarks. The primary
contributions of this work are three-fold:

\begin{enumerate}
    \item \textbf{Automated Multi-Sensor Annotation Pipeline:} We designed a
          structured pipeline that synchronizes four visual sensor modalities (RGB,
          Infrared, Event, and Marigold Depth) via Dynamic Time Warping (DTW),
          cross-correlation, and optical flow. Using a three-stage LLM prompt chain (Captioning $\rightarrow$
          Question Generation $\rightarrow$ Answering) combined with fine-grained
          atomic QA splitting, the pipeline converts raw multi-sensor recordings into
          structured QA candidate pools.

    \item \textbf{Two-Tier Quality Control and Late Fusion Framework:} We introduced
          a rigorous quality assurance mechanism combining rule-based sanitization
          and LLM-assisted semantic verification across five quality dimensions and
          hallucination risk estimation. This is coupled with a late fusion module
          governed by category-driven reliability profiles and relation-aware
          cross-modal support scoring. From 2,730 initial raw blocks (yielding 7,310
          atomic candidates), the framework filters out malformed or ungrounded entries
          and retains 5,465 high-quality visual QA items for the benchmark.

    \item \textbf{Controlled VLM Benchmarking and Empirical Insights:} We established
          three standardized evaluation protocols—Source-Modality Baseline, Primary
          Modality Replacement (21,324 evaluations across 5,331 paired questions), and
          Native Video Adapter. Empirical evaluation across 4B and 8B parameter
          variants of Qwen3-VL, InternVL3.5, and Molmo2 revealed that model architecture
          and evidence representation play a more decisive role in modality robustness
          than parameter scale alone. Notably, Qwen3-VL achieved superior performance
          across all visual modalities (reaching 57.49\% strict accuracy at 8B scale).
\end{enumerate}

While these contributions demonstrate the viability of automated multi-sensor
annotation and controlled VLM benchmarking, the framework remains subject to certain
scope caveats—such as evaluator LLM bias, single-modality grounding scope, and
cross-sensor imbalances—which are discussed in detail in Chapter~\ref{ch:evaluation}.


\section{Future Directions}
\label{sec:conclusion_future_work}

Building upon the findings of this project, several promising avenues exist for future research:

\begin{itemize}
    \item \textbf{Human-in-the-Loop Validation:} Conducting a stratified human
          validation study across retained, review-recommended, and rejected QA
          items to quantify agreement between human annotators and the LLM-assisted
          semantic verifier.

    \item \textbf{Joint Multimodal Fusion Protocols:} Extending the evaluation suite to
          incorporate dedicated multi-input protocols that evaluate models on
          cross-modal disambiguation and joint reasoning tasks across simultaneous
          multi-sensor inputs.

    \item \textbf{Scaling to Larger-Scale Multi-Scene Datasets:} Extending the
          automated annotation pipeline and benchmarking framework to larger,
          more diverse egocentric and multi-sensor video datasets to assess
          scalability, generalize reliability profiles, and calibrate cross-modal
          support thresholds across varied real-world environments.

    \item \textbf{Enhancing QA Complexity and Reasoning Depth:} Expanding from
          atomic, short-form QA pairs toward multi-step spatial-temporal reasoning
          and descriptive open-ended questions to challenge next-generation VLMs
          beyond surface-level visual recognition.
\end{itemize}
